\documentclass[11pt,a4paper]{article}

% Packages
\usepackage[left=1.5cm,right=1.5cm,top=1.5cm,bottom=1.5cm]{geometry}
\usepackage{fontspec}
\usepackage{xcolor}
\usepackage{hyperref}
\usepackage{titlesec}
\usepackage{enumitem}
\usepackage{tabularx}
\usepackage{newunicodechar}

% Farben
\definecolor{darkblue}{RGB}{0,47,108}
\definecolor{lightgray}{RGB}{100,100,100}

% Font
\setmainfont{Helvetica Neue}
\setsansfont{Helvetica Neue}
\renewcommand{\familydefault}{\sfdefault}

% Hyperref Setup
\hypersetup{
    colorlinks=true,
    linkcolor=darkblue,
    urlcolor=darkblue,
    pdfauthor={Keanu Fuchs},
    pdftitle={Lebenslauf Keanu Fuchs}
}

% Section Formatting
\titleformat{\section}
    {\large\bfseries\color{darkblue}}
    {}
    {0em}
    {}
    [\titlerule]

\titlespacing*{\section}{0pt}{8pt}{4pt}

% Custom Commands
\newcommand{\cventry}[3]{
    \noindent
    \begin{minipage}[t]{0.75\textwidth}
        \textbf{#1}\\
        \textit{\small #2}
    \end{minipage}
    \hfill
    \begin{minipage}[t]{0.23\textwidth}
        \raggedleft\textcolor{lightgray}{\small #3}
    \end{minipage}
    \par\vspace{6pt}
}

\newcommand{\cveducation}[3]{
    \noindent
    \begin{minipage}[t]{0.75\textwidth}
        \textbf{#1}\\
        \textit{\small #2}
    \end{minipage}
    \hfill
    \begin{minipage}[t]{0.23\textwidth}
        \raggedleft\textcolor{lightgray}{\small #3}
    \end{minipage}
    \par\vspace{6pt}
}

% List Settings
\setlist[itemize]{leftmargin=*,topsep=0pt,itemsep=1pt,parsep=0pt}

% Document
\begin{document}
\pagestyle{empty}

% Header
\noindent
\begin{minipage}[t]{0.65\textwidth}
    {\huge\bfseries\color{darkblue}Keanu Fuchs}\\[2pt]
    {\large\color{lightgray}System Engineer | Network Automation}\\[6pt]
    \textcolor{lightgray}{
        Schlesienstraße 5, 53119 Bonn\\
        \href{tel:+491746352391}{+49 174 6352391} | \href{mailto:kn.fchs@gmail.com}{kn.fchs@gmail.com}\\
        \href{https://linkedin.com/in/keanufuchs}{linkedin.com/in/keanufuchs} |
        \href{https://github.com/keanufuchs}{github.com/keanufuchs}
    }
\end{minipage}

\vspace{10pt}

% Profil
\section*{Profil}
System Engineer und Dualer Student mit Fokus auf Netzwerkautomatisierung (Cisco) und Softwareentwicklung. Spezialisiert auf skalierbare Provisionierung mittels Python sowie die Integration moderner KI-Ansätze in Network Operations.

\vspace{4pt}

% Kompetenzen
\section*{Kompetenzen}
\begin{tabularx}{\textwidth}{@{}X@{\hspace{2em}}X@{\hspace{2em}}X@{}}
    \textbf{Network Automation} & \textbf{Development} & \textbf{APIs \& Protokolle}\\
    Cisco Catalyst Center, PnP, & Python, Ansible, Flask, & REST, NETCONF/YANG,\\
    Day-0/Day-N, Jinja2 & PyATS, Git, GitLab CI/CD & Cisco APIs, Meraki\\
\end{tabularx}

\vspace{4pt}

% Berufserfahrung
\section*{Berufserfahrung}

\cventry{Dualer Student Informatik (B.Sc.) \& System Engineer}{Bechtle IT-Systemhaus Bonn/Köln}{Sep. 2023 – heute}
\begin{itemize}
    \item Netzwerkautomatisierung im Cisco-Enterprise-Umfeld (Catalyst Center, Python, REST, NETCONF/YANG)
    \item Automatisierung von Provisionierung \& Lifecycle-Prozessen (Day-0/Day-N)
    \item API-Integrationen, CI/CD-Pipelines (GitLab) und Docker-Containerisierung
    \item Technischer Lead bei Enterprise-Rollouts (Logistik- und Industriestandorte)
\end{itemize}
\vspace{4pt}

\cventry{Junior System Engineer}{Bechtle IT-Systemhaus Bonn/Köln}{Jun. 2023 – Sep. 2023}
\begin{itemize}
    \item Cisco Catalyst Center: Provisionierung, Plug-and-Play, Day-0/Day-N Templates (Jinja2)
\end{itemize}
\vspace{4pt}

\cventry{Fachinformatiker für Systemintegration (Ausbildung)}{Bechtle IT-Systemhaus Bonn/Köln}{Sep. 2021 – Jun. 2023}
\begin{itemize}
    \item Cisco Networking, Enterprise-Netzwerkumgebungen, Abschlussprojekt: Automatisierter Geräteaustausch über Catalyst Center
\end{itemize}

\vspace{4pt}

% Ausbildung
\section*{Ausbildung}

\cveducation{B.Sc. Informatik (dual) – Schwerpunkt AI \& Network Automation}{Rheinische Hochschule Köln}{Sep. 2023 – heute}
\cveducation{Fachinformatiker für Systemintegration (IHK)}{Heinrich-Hertz-Europakolleg Bonn}{Sep. 2021 – Jun. 2023}
\cveducation{Allgemeine Hochschulreife (Abitur)}{Städtisches Gymnasium Rheinbach}{Jun. 2021}

\vspace{4pt}

% Zertifikate & Sprachen
\section*{Zertifikate \& Sprachen}
\noindent
\textbf{Cisco CCNA} \hfill \textcolor{lightgray}{\small Feb. 2024 – Feb. 2027}\\[4pt]
\textbf{Deutsch:} Muttersprache \quad \textbf{Englisch:} Sehr gut \quad \textbf{Französisch:} Grundkenntnisse

\end{document}
